% ---
% id: meat-core/req:split-form-gradient
% kind: req
% title: "Split-form evaluation of the trace-power gradient"
% summary: >-
%   The engine evaluates the gradient of the pinched trace power in the split
%   form of [[thm:gradient]] (i): divided differences over the positive
%   spectrum, pinched operators on the ambient space. It compresses the whole
%   computation only to a pinching-invariant space containing the branch support.
% prominence: standard
% status: implemented
% uses: [thm:gradient, def:branch-support]
% mentions: [lem:branch-support]
% assumes: []
% refs: []
% provenance: newly-derived
% implements: certiqs_key.meat_core.objective.gradient:pullback_trace_power_gradient
% lean:
% ---

A branch operator is in general singular, and its
\tref{maximal branch support} is in general not invariant under the
\tref{key pinching} (Lemma~\ref{lem:branch-support}\,(iv)); compressing the whole
computation to its support then computes a different quantity, as the
rank-one witness in the proof of Theorem~\ref{thm:gradient} shows.

\begin{requirement}\label{req:split-form-gradient}
The engine evaluates the gradient of the pinched trace power at a branch
operator $X$ in the split form of Theorem~\ref{thm:gradient}\,(i): the
divided-difference map $DX^{p}[\cdot]$ compressed to
$\operatorname{supp}X$, its kernel summed over the strictly positive
eigenvalues of $X$ only, and $\cZ(X^{p})$, $Y^{\alpha-1}$ and
$\cZ(Y^{\alpha-1})$ formed on the ambient branch space. It compresses the
whole computation to a subspace $S'$ only where $S'$ contains the maximal
branch support $W_{\hat c,s}$ of Definition~\ref{def:branch-support} and
$Z_jS'\subseteq S'$ for every $j$.
\end{requirement}
